\section{Related Work}
\label{sec:related}

SE belongs to a broader family of graph-complexity, graph-clustering, and
graph-compression methods, but it differs from most classical objectives in
treating hierarchical organization as the primary object to be decoded.
The closest reference points for understanding its role are modularity-based
community detection, flow-based compression, spectral partitioning, and modern
deep graph clustering.
These neighboring lines of work provide useful comparisons because they often
pursue similar structural goals while relying on different assumptions about
what counts as a good partition.

\paragraph{Modularity-based community detection.}
Modularity compares observed within-community edge density against a
configuration-model null graph~\cite{newman2004findingmodularity}.
Louvain and Leiden-style optimization make this objective highly scalable, and
Leiden further improves partition connectivity and refinement
stability~\cite{traag2019leiden}.
These methods remain strong default baselines for topology-driven clustering,
but they are fundamentally flat and are known to exhibit resolution effects.
SE differs in that it evaluates the coding cost induced by an encoding tree,
so the partition is interpreted as an information hierarchy rather than only a
density surplus.

\paragraph{Flow-based compression.}
Flow-based compression methods such as the map equation and Infomap encode
random-walk trajectories using module codebooks~\cite{rosvall2008maps}.
They are conceptually close to SE because both explain graph organization
through compression, but the object being compressed is different.
Infomap compresses flows generated by random walks, whereas SE compresses
structural uncertainty under graph volumes, cuts, and tree refinement.
This distinction becomes important when the graph is interpreted as an
information system rather than an explicit navigation process.

\paragraph{Spectral and cut-based methods.}
Spectral and cut-based methods use eigenvectors of graph Laplacians to expose
low-conductance partitions~\cite{shi2000normalized,ng2002spectral}.
They provide powerful relaxations for balanced graph partitioning and remain
important for initialization, diagnostics, and theoretical comparison.
In most forms, however, spectral objectives are flat and externally
parameterized by the target number of clusters.
SE can connect to spectral information, including the literature around
von Neumann entropy, but its target remains a minimum-description hierarchy
rather than an eigenspace partition alone.

\paragraph{Deep graph clustering and pooling.}
Deep graph clustering and pooling methods such as DiffPool, MinCutPool, and
DMoN learn soft assignments inside graph neural network architectures~\cite{ying2018diffpool,bianchi2020mincutpool,tsitsulin2020dmon}.
They are effective when node attributes and task supervision provide extra
semantic signal, but their objectives are often architectural regularizers
rather than standalone measures of structural organization.
Recent SE-based neural methods, including LSENet, SEP-style pooling, and
DeSE, help bridge this gap by making structural criteria usable inside learned
models~\cite{wu2022pooling,sun2024lsenet,zhang2025dese}.
One goal of the present survey is to place these neural methods on the same
conceptual map as the original combinatorial SE algorithms and their
domain-specific extensions.

Taken together, these neighboring paradigms clarify what is distinctive about
SE.
It is not simply another flat graph-clustering score, nor only a regularizer
for graph neural networks.
Its central claim is that graph organization should be understood through the
coding efficiency of hierarchical structure, and that this perspective can be
used both analytically and algorithmically across a wide range of graph
intelligence tasks.
